\chapter{The Gaussian Chain by Linear Algebra}\label{ch:gaussian}

The stationary Gaussian AR(1) chain is the one model in this thesis whose joint
score can be written down in closed form and checked by hand. This chapter does
that with matrices: the clean law, its dense covariance and sparse precision,
the score in two equivalent forms, and the lifecycle the coupling structure
passes through as the diffusion time grows. Chapter~\ref{ch:gaussian-bp} then
recomputes the same object by message passing and shows the two agree.
Everything here is exact unless marked as a numerical check, and each boxed
formula is covered by a named check of the audit suite
(Appendix~\ref{app:repro}).
\section{The joint law and the object of study}\label{sec:g-model}

A data point is a trajectory $a = (a_0, \dots, a_{L-1})$ of the stationary
AR(1) chain \eqref{eq:ar1} with $\sigma_\eta^2 = 1 - \alpha^2$. Each frame
passes independently through the OU channel \eqref{eq:ou-channel} for a
time $t$,
\begin{equation}
  x_k \;=\; e^{-t} a_k + \sqrt{\Deltat}\;\xi_k,
  \qquad \xi_k \sim \Nd(0,1)\ \text{i.i.d.},
  \label{eq:g-channel}
\end{equation}
and the object of study is the joint score of the noisy sequence,
\begin{equation}
  \score(x, t) \;=\; \nabla_x \log P_t(x) \;\in\; \R^L,
  \qquad
  P_t(x) = \int_{\R^L} P_0(a) \prod_{k=0}^{L-1}
    \Nd\big(x_k;\, e^{-t}a_k,\, \Deltat\big)\, \dd a .
  \label{eq:g-object}
\end{equation}
The gradient is taken with respect to all $L$ noisy frames simultaneously,
under the law of the whole noisy trajectory
(Section~\ref{sec:tweedie}).

\section{The clean law: dense covariance, sparse precision}
\label{sec:g-clean}

The chain factorisation \eqref{eq:chain-factorisation} gives the clean
joint density
\begin{equation}
  P_0(a)
  = \Nd(a_0; 0, 1) \prod_{k=0}^{L-2}
    \Nd\big(a_{k+1};\, \alpha a_k,\, \sigma_\eta^2\big)
  = \Nd(a;\, 0,\, \Sigz).
  \label{eq:g-P0}
\end{equation}

\begin{proposition}[Stationary covariance]\label{prop:g-Sigma0}
$(\Sigz)_{ij} = \alpha^{|i-j|}$: a dense, symmetric Toeplitz matrix with
geometrically decaying off-diagonals, $\gamma(d) = \alpha^{|d|}$.
\end{proposition}

\begin{proposition}[Tridiagonal precision]\label{prop:g-Q0}
$\Qz := \Sigz^{-1}$ is exactly tridiagonal:
\begin{equation}
  (\Qz)_{kk} =
  \begin{cases}
    \dfrac{1}{\sigma_\eta^2}, & k \in \{0, L-1\},\\[8pt]
    \dfrac{1+\alpha^2}{\sigma_\eta^2}, & 0 < k < L-1,
  \end{cases}
  \qquad
  (\Qz)_{k, k\pm 1} = -\frac{\alpha}{\sigma_\eta^2},
  \label{eq:g-Q0}
\end{equation}
and every entry at distance $\geq 2$ vanishes identically.
\end{proposition}

\begin{toolbox}{Deriving the tridiagonal precision from the log-density}
For a centred Gaussian, the precision is the Hessian of the negative
log-density. From \eqref{eq:g-P0},
\begin{equation*}
  -2 \log P_0(a)
  = \frac{a_0^2}{\sigma_0^2}
  + \sum_{k=0}^{L-2} \frac{(a_{k+1} - \alpha a_k)^2}{\sigma_\eta^2}
  + \text{const},
  \qquad \sigma_0^2 = 1 .
\end{equation*}
Read off the quadratic form term by term. The only term coupling $a_k$ and
$a_{k+1}$ is the transition $(a_{k+1} - \alpha a_k)^2/\sigma_\eta^2$, with
cross coefficient $-2\alpha/\sigma_\eta^2$, hence precision entry
$-\alpha/\sigma_\eta^2$; no term couples variables at distance $\geq 2$,
so those entries are exactly zero. An interior $a_k$ appears in its left
transition with coefficient $1/\sigma_\eta^2$ and in its right transition
with $\alpha^2/\sigma_\eta^2$, summing to $(1+\alpha^2)/\sigma_\eta^2$. At
the left boundary the prior term $1/\sigma_0^2 = (1-\alpha^2)/\sigma_\eta^2$
plus the right transition $\alpha^2/\sigma_\eta^2$ give $1/\sigma_\eta^2$;
the right boundary is symmetric. \hfill$\square$
\end{toolbox}

\begin{remark}[Tridiagonality is Markovianity, not Gaussianity]
\label{rmk:g-markov-sparsity}
A zero in the precision matrix is a conditional independence given all
other variables; by the Hammersley--Clifford theorem
(Section~\ref{sec:mrf}), factorisation over the cliques of a graph and
conditional independence with respect to that graph are the same property.
Gaussianity entered the derivation only to identify the precision with the
Hessian of $-\log P_0$; the support of that Hessian, nearest neighbours
only, comes from the chain factorisation alone. Any Markov chain with
smooth transitions has this sparsity pattern; the Gaussian case makes the
entries constants.
\end{remark}

The covariance is dense (every frame correlates with every other) while
the precision is tridiagonal: correlation is global, conditional
dependence is local. The score is governed by the second object.

\section{The score, two ways}\label{sec:g-score}

\subsection{Linear algebra}

\begin{proposition}[Noised law]\label{prop:g-Sigt}
$P_t(x) = \Nd(x;\, 0,\, \Sigt)$ with
\begin{equation}
  \boxed{\;\Sigt \;=\; e^{-2t}\, \Sigz \;+\; \Deltat\, I_L\;}
  \label{eq:g-Sigt}
\end{equation}
\end{proposition}
This is \eqref{eq:ou-cov-map}. The diagonal is
$e^{-2t} + \Deltat = 1$ at every $t$ (variance preservation); only the
off-diagonal entries, $\alpha^{|i-j|} e^{-2t}$, depend on the diffusion
time, because the noise of different frames is independent. Since
$\log P_t(x) = -\tfrac{1}{2} x^\T \Sigt^{-1} x + \text{const}$,
\begin{equation}
  \boxed{\;\score(x, t) \;=\; -\,\Qt\, x,
  \qquad \Qt := \Sigt^{-1}.\;}
  \label{eq:g-score-matrix}
\end{equation}
The joint score is linear in $x$, the hallmark of Gaussianity, and its
entire structure is the noisy precision matrix $\Qt$: which frames couple,
how strongly, and at which $t$, are questions about this one matrix,
answered in Sections~\ref{sec:g-lifecycle} and~\ref{sec:g-bulk}.

\subsection{Bayesian inference}

The score--posterior identity \eqref{eq:tweedie-joint} expresses the same
score through the posterior mean of the clean chain,
\begin{equation}
  \score_k(x, t)
  = \frac{e^{-t}\, \E[a_k \,|\, x] - x_k}{\Deltat} ,
  \label{eq:g-tweedie}
\end{equation}
and the posterior is Gaussian with parameters used throughout:

\begin{proposition}[Posterior of the clean chain]\label{prop:g-posterior}
$P(a \,|\, x) = \Nd(a;\, m,\, J^{-1})$ with
\begin{equation}
  \boxed{\;
  J \;=\; \frac{e^{-2t}}{\Deltat}\, I \;+\; \Qz ,
  \qquad
  J\, m \;=\; h := \frac{e^{-t}}{\Deltat}\, x .
  \;}
  \label{eq:g-posterior}
\end{equation}
\end{proposition}

\begin{toolbox}{Posterior precision and field, by completing the square}
Bayes' rule with the likelihood \eqref{eq:g-channel}:
\begin{align*}
  -2 \log P(a | x)
  &= a^\T \Qz\, a
  + \sum_{k} \frac{(x_k - e^{-t} a_k)^2}{\Deltat}
  + \text{const} \\
  &= a^\T \Big( \Qz + \tfrac{e^{-2t}}{\Deltat} I \Big) a
  - 2\, \tfrac{e^{-t}}{\Deltat}\, x^\T a
  + \text{const}.
\end{align*}
Matching against the generic form $a^\T J a - 2 h^\T a$ gives
\eqref{eq:g-posterior}. Two structural observations. (i) The likelihood is
diagonal in $a$: observing each frame through its own channel adds the
channel precision $e^{-2t}/\Deltat$ to every diagonal entry of the prior
precision and touches nothing else. (ii) Consequently $J$ is tridiagonal
at every $t$: the \emph{posterior} of the clean chain remains a
nearest-neighbour Markov field; what fills up with $t$ is the precision
$\Qt$ of the noisy marginal law. Keeping $J$ (tridiagonal, posterior) and
$\Qt$ (dense, noisy marginal) distinct prevents much confusion.
\hfill$\square$
\end{toolbox}

Substituting $m = J^{-1} h$ into \eqref{eq:g-tweedie} must return
$-\Qt x$; the operator identity behind this is Woodbury's
(Appendix~\ref{app:gauss}), and the audit suite checks the two routes
agree to $2.7 \times 10^{-15}$ across sweeps of $(L, \alpha, t)$. The
agreement is an end-to-end check of two computations built from different
formulas.

\subsection{The smallest example, $L = 2$}\label{sec:g-K2}

For $L = 2$, with $r := \alpha\, e^{-2t}$,
\begin{equation}
  \Sigt = \begin{pmatrix} 1 & r \\ r & 1 \end{pmatrix}
  \qquad\Longrightarrow\qquad
  \score_0(x, t) = -\,\frac{x_0 - r\, x_1}{1 - r^2},
  \qquad
  \score_1(x, t) = -\,\frac{x_1 - r\, x_0}{1 - r^2}.
  \label{eq:g-K2-score}
\end{equation}
By contrast, each $x_k$ is marginally $\Nd(0,1)$ at every $t$ (variance
preservation), so the per-frame marginal score is $-x_k$: independent of
$t$ and of the other frame, it carries no information about the dynamics.
The joint score contains the cross term $-r x_1/(1-r^2)$: the denoiser of
frame $0$ uses evidence from frame $1$ with a weight that equals $\alpha$
at $t = 0$ and decays as $e^{-2t}$. This is the coupling structure whose
full lifecycle the next section derives, in miniature.

\section{The spectral form and the lifecycle of the precision}
\label{sec:g-lifecycle}

Between the tridiagonal $\Qz$ at $t = 0$ and the identity at $t = \infty$ lies
the whole lifecycle of the coupling structure. The analysis rests on the
spectral theorem
and its affine-family corollary: affine
images of a symmetric matrix share its eigenbasis,
$(cA + dI)\, u_i = (c\,\omega_i + d)\, u_i$. Applying this to the
forward map \eqref{eq:g-Sigt}, which is affine in $\Sigz$, gives the
representation from which everything in this section follows.

\begin{theorem}[Spectral form]\label{thm:g-three-forms}
Let $\Sigz = U \operatorname{diag}(\omega_i) U^\T$. Then for every
$t \geq 0$,
\begin{equation}
  \boxed{\;
  \Sigt = U \operatorname{diag}\big( \lambda_i(t) \big) U^\T,
  \qquad
  \Qt = U \operatorname{diag}\!\Big( \frac{1}{\lambda_i(t)} \Big) U^\T,
  \qquad
  \lambda_i(t) = e^{-2t} \omega_i + \Deltat .
  \;}
  \label{eq:g-Qt-spec}
\end{equation}
\end{theorem}

Diffusion never rotates the problem: the eigenvectors of $\Sigt$ and
$\Qt$ are those of $\Sigz$, frozen for all $t$, while each eigenvalue
flows along the affine path $\omega_i \mapsto e^{-2t}\omega_i + \Deltat$,
monotonically towards $1$. Because $\Sigz$ is symmetric Toeplitz, its
eigenvectors are asymptotically the Fourier modes of the sequence
: in frequency space the denoising problem is $L$
independent scalar problems at every diffusion time. Two computational
corollaries are recorded for later use: factoring $\Deltat$ out of
\eqref{eq:g-Sigt} gives the SNR form
\begin{equation}
  \Qt = \tfrac{1}{\Deltat} \big( I + \gamma_t\, \Sigz \big)^{-1},
  \qquad \gamma_t := e^{-2t}/\Deltat ,
  \label{eq:g-Qt-snr}
\end{equation}
and multiplying \eqref{eq:g-Sigt} by $\Qz$ and inverting gives the
resolvent form
\begin{equation}
  \Qt = e^{2t}\, \Qz \big( I + c_t \Qz \big)^{-1}
      = e^{2t} \sum_{m \geq 0} (-c_t)^m\, \Qz^{\,m+1},
  \qquad c_t := e^{2t} - 1 ,
  \label{eq:g-Qt-res}
\end{equation}
whose Neumann series converges for $\|c_t \Qz\| < 1$, i.e.\ at small $t$.

\subsection{Losing tridiagonality: when, how, and how much}
\label{sec:g-tridiag-loss}

Sparsity is a property of a matrix in a basis. In the eigenbasis the
problem is diagonal at all times; in the frame basis, the operationally
meaningful one, the zeros of $\Qz$ are a cancellation that diffusion
destroys. The destruction is quantifiable, and it answers three questions:
when tridiagonality is lost, how, and how much structure survives at
each $t$.

\paragraph{When: immediately.} By \eqref{eq:g-Qt-spec} every entry of
$\Qt$ is analytic in $t$. The distance-$d$ entries, identically zero at
$t = 0$ for $d \geq 2$, are non-zero for every $t > 0$: no exact zero
survives any positive diffusion time.

\paragraph{How: one diagonal per order in $t$.} The loss is nevertheless
graded, and the resolvent form gives the exact rate at which each diagonal
switches on.

\begin{theorem}[Band-fill leading order]\label{thm:g-bandfill}
For small $t$ and frame distance $d \geq 1$,
\begin{equation}
  \boxed{\;
  (\Qt)_{i, i+d}
  \;=\; (-1)^{d-1}\, (2t)^{d-1}\, \big( \Qz^{\,d} \big)_{i, i+d}
  \;+\; O(t^{d}) .
  \;}
  \label{eq:g-bandfill}
\end{equation}
\end{theorem}

\begin{toolbox}{Band fill from the Neumann series}
In the resolvent form \eqref{eq:g-Qt-res}, substitute
$c_t = 2t + 2t^2 + O(t^3)$ and $e^{2t} = 1 + O(t)$:
\begin{equation*}
  \Qt = \big(1 + O(t)\big)
  \Big[ \Qz - c_t \Qz^2 + c_t^2 \Qz^3 - \cdots \Big].
\end{equation*}
$\Qz$ is tridiagonal, so $\Qz^{\,m}$ has bandwidth $m$: the entry at
distance $d$ receives its first contribution from the single term $m = d$,
with prefactor $(-c_t)^{d-1} = (-1)^{d-1}(2t)^{d-1}(1 + O(t))$; every
later term is $O(t^d)$. Each additional frame of coupling range costs one
power of $t$: the band fills outward, one diagonal per order.
\hfill$\square$
\end{toolbox}

\begin{remark}[The prefactor carries no factorial, and why that is checked]
\label{rmk:g-correction}
Equation~\eqref{eq:g-bandfill} has no $1/(d-1)!$: the leading distance-$d$
term is a single Neumann term, with no combinatorial sum behind it. The
distinction is easy to miss and worth verifying numerically, because slope
fits in $\log t$ are blind to constant prefactors, a wrong constant
survives every check on the exponent. The audit therefore tests the
coefficient itself, and it separates the two candidates cleanly: the
no-factorial form matches the measured $(\Qt)_{i,i+d}/t^{d-1}$ to better
than $0.9\%$ for every $d \in \{1,\dots,5\}$ (at $\alpha = 0.9$, $L = 25$),
while a factorial version would be wrong by $86\%$ already at $d = 3$. Every
prefactor in this thesis is audited on the same grounds, not only every
exponent.
\end{remark}

\begin{figure}[t]\centering
\includegraphics[width=0.94\linewidth]{fig_band_fill.pdf}
\caption[The band-fill law for the posterior precision]{Band-fill law. Measured $|(\Qt)_{i,i+d}|$ against $t$ (dots) and
the leading order \eqref{eq:g-bandfill} (solid lines, no fitted
parameters), for $\alpha = 0.9$, $L = 25$, $d = 1, \dots, 5$. The slopes
are $d - 1$; agreement at small $t$ is at the percent level for every
$d$.}
\label{fig:g-bandfill}
\end{figure}

\paragraph{How much: the off-tridiagonal mass.}
Define the off-tridiagonal Frobenius mass of $\Qt$: the norm of everything
outside the band that Markov structure would allow. At small $t$ the
band-fill law makes it rise from zero, dominated by the $d = 2$ diagonal,
linearly in $t$. At large $t$, expanding the inverse of
$\Sigt = I + e^{-2t}(\Sigz - I)$ gives
\begin{equation}
  \Qt = I - e^{-2t}(\Sigz - I) + e^{-4t}(\Sigz - I)^2 + O(e^{-6t}),
  \label{eq:g-larget}
\end{equation}
so $\| \Qt - I \|_F \simeq e^{-2t} \|\Sigz - I\|_F$: all structure,
off-tridiagonal and tridiagonal alike, decays at the rate $e^{-2t}$, and
the score approaches the isotropic $-x$. In between, the off-tridiagonal
mass peaks (Figure~\ref{fig:g-tridiagloss}): the noisy sequence is most
non-Markov at intermediate diffusion times. The lifecycle in one line:
tridiagonal at $t = 0$; densifying at rate $t^{d-1}$ per diagonal $d$;
maximally coupled at intermediate $t$; isotropic as $t \to \infty$ at rate
$e^{-2t}$.

\begin{figure}[t]\centering
\includegraphics[width=0.94\linewidth]{fig_tridiag_loss.pdf}
\caption[Loss and recovery of tridiagonality with diffusion time]{Structure across diffusion time ($\alpha = 0.8$, $L = 40$).
$\|\Sigt - \Sigz\|_F$ (navy) grows monotonically as the covariance is
driven to the identity; the off-tridiagonal Frobenius mass of $\Qt$ (red)
rises from zero, peaks at intermediate $t$, and decays to zero: the band
fills, then empties.}
\label{fig:g-tridiagloss}
\end{figure}
